\documentclass[11pt]{article}
\usepackage[T1]{fontenc}
\usepackage[margin=1in]{geometry}
\usepackage{amsmath,amssymb,amsthm,mathtools,microtype,hyperref}
\hypersetup{hidelinks,pdftitle={Fermat control for abelian covers of arbitrary exponent}}
\newtheorem{theorem}{Theorem}[section]
\newtheorem{lemma}[theorem]{Lemma}
\newtheorem{proposition}[theorem]{Proposition}
\newtheorem{corollary}[theorem]{Corollary}
\theoremstyle{remark}\newtheorem{remark}[theorem]{Remark}
\newcommand{\Q}{\mathbf Q}\newcommand{\C}{\mathbf C}
\newcommand{\Z}{\mathbf Z}\newcommand{\PP}{\mathbf P}
\DeclareMathOperator{\SL}{SL}\DeclareMathOperator{\Sp}{Sp}
\DeclareMathOperator{\PGL}{PGL}\DeclareMathOperator{\tr}{tr}
\DeclareMathOperator{\Hom}{Hom}\DeclareMathOperator{\Gal}{Gal}
\DeclareMathOperator{\HC}{HC}\DeclareMathOperator{\CH}{CH}
\DeclareMathOperator{\Sym}{Sym}\DeclareMathOperator{\im}{im}
\title{Fermat control of Hodge classes for abelian covers\\of arbitrary exponent}
\author{Research continuation for \texttt{abelian-cover-hodge-lean}}
\date{5 September 2026 --- written proof draft, not independently reviewed}
\begin{document}
\maketitle
\begin{abstract}
We give a proposed extension of the repository's odd-exponent reduction to
arbitrary finite abelian deck groups. The extra arguments are an even-order
rank-two graph test and a uniform three-point model for every all-definite
primitive factor. The high-rank parity argument already occurs in
Menet--Nguyen and is explicitly credited. With the classical primitive
Fermat-determinant construction, these arguments give an algebraic-correspondence
presentation of the connected-monodromy invariants on every power of a
very general covering Jacobian. Consequently the all-powers rational Hodge
conjecture for the degree-$N$ Fermat Jacobian is equivalent to its validity
for every full abelian-cover family of exponent dividing $N$. The reverse
implication simply includes the Fermat curve itself. This is a proposed
proof with stated classical inputs, not a completed Lean proof or a
certification of priority or journal significance.
\end{abstract}

\paragraph{Checkpoint and scope.}
This note continues the repository checkpoint\\
{\small\texttt{add067109eb00dc9bd3e4cd3dcdcc2852067dd87}},
without modifying its frozen sources. The inherited correspondence and rational
assembly arguments are in \cite[Sections 4--6]{General} and
\cite[Section 7]{Structural}. We restate their precise uses below and supply
the new parity and constant-factor arguments. All varieties and Hodge
structures are over $\C$ and $\Q$, respectively. No statement about all
special fibers, integral Hodge classes, nonabelian covers, or restricted
branch subfamilies is made.

\section{Statements}
Let $N\geq2$, let $G$ be a finite abelian group of exponent dividing $N$,
and let $a_1,\ldots,a_r\in G\setminus\{0\}$ generate $G$, with
$\sum_i a_i=0$ and $r\geq3$. Let $B=M_{0,r}$ be the full ordered branch
configuration space. Passing to a finite cover to choose the family does
not change the connected monodromy or the meaning of very general. Write
$X_b\to\PP^1$ for its connected $G$-cover and $A_b=J(X_b)$.
Let $F_N^j\subset\PP^{j+1}$ be a smooth degree-$N$ Fermat hypersurface;
$F_N=F_N^1$ is the Fermat curve.

\begin{theorem}[Fermat-source presentation]\label{thm:presentation}
For every power $q$ and degree $k$, the connected-monodromy invariant part
$H^k(A_b^q,\Q)^M$ is the image of a finite direct sum of maps induced by
rational algebraic correspondences. Their sources are tensor products of
primitive Fermat-determinant summands, first cohomology of isogeny factors
of $J(F_N)$, and effective Tate factors. Each determinant summand is an
algebraic direct summand of some $H^h(F_d^h,\Q)$ with $d\mid N$.
The maps present the entire invariant space, including mixed-conductor
classes and finite-monodromy factors.
\end{theorem}
This is a cohomological presentation by algebraic maps. It does not assert
an algebraic projector onto the full invariant space or an algebraic
splitting of an arbitrary Hodge surjection.

\begin{theorem}[All-exponent transfer]\label{thm:transfer}
Assume the rational Hodge conjecture for every power of $J(F_N)$. Then,
for every datum above, it holds for every power of $A_b$ at every point
outside one countable union of proper closed algebraic subsets of $B$.
In particular the following assertions are equivalent:
\[
 \HC(J(F_N)^q)\text{ for every }q
 \quad\Longleftrightarrow\quad
 \begin{gathered}
 \HC(A_b^q)\text{ for every }q\text{ and every such datum,}\\
 \text{with }b\text{ very general in its full branch family.}
 \end{gathered}
\]
The quantification on the right includes the rigid three-branch families.
\end{theorem}

\begin{corollary}[The universal fixed-part Hodge category]\label{cor:category}
Let $\mathcal T_N$ be the smallest rigid tensor subcategory of polarizable
rational Hodge structures containing Tate objects and all spaces
$H^k(A_b^q,\Q)^M$ above, as the full branch datum varies and $b$ is very
general. Then
\[
 \mathcal T_N=\langle H^1(F_N,\Q),\Q(1)\rangle.
\]
This is an equality of Hodge-theoretic tensor categories, not an assertion
that the corresponding Chow-motive categories have been identified.
\end{corollary}

\begin{corollary}[Established Fermat degrees as input]\label{cor:aoki}
Subject to the proposed proof of Theorem~\ref{thm:transfer}, its conclusion
is unconditional for
\[
 N=p^e\text{ or }2p^e\quad(p\text{ prime},\ e\geq1),
\]
and for
\[
 N=2^a3^b5^c7^d\geq2,
 \qquad a,b,c,d\geq0,\qquad c=0\text{ or }d=0.
\]
Here the cycle input is Aoki's theorem, not a new Fermat-cycle result.
\end{corollary}
The first list is also restated in \cite[Theorem 1.2]{Aoki02}, together
with the second. In particular, the statement includes all finite abelian
$2$-groups and many even mixed-exponent groups. It does not infer the
Fermat premise at unrestricted $N$.

\section{Primitive factors and a dichotomy without conductor exceptions}
Fix $\zeta_N$. An evaluated character is a word
$c=(c(a_1),\ldots,c(a_r))\in(\Z/N)^r$. Its order is denoted by $d(c)$.
If $d=d(c)$, write $c_i=(N/d)b_i$, with $b_i\in\Z/d$.
The active $b_i$ are nonzero, sum to zero, and generate $\Z/d$.
Write $s$ for their number and $h=s-2$.

One representative of each unit orbit of characters gives the rational
primitive decomposition of $H^1(X_b,\Q)$. More precisely, for the degree-$d$
cyclic quotient $Y_c=X_b/\ker c$, take the primitive factor
\[
 U_c=\ker\Phi_d(T_c^*)\subset H^1(Y_c,\Q),
 \qquad
 U_c\otimes\C=\bigoplus_{u\in(\Z/d)^\times}V_{uc}.
\]
It has rank $h$ over $K_d=\Q(\zeta_d)$. Quotient pullback, norm, and
cyclotomic idempotents are rational algebraic maps. One must not sum the
whole cohomology of every cyclic quotient: that double counts the
lower-conductor factors. For example, the primitive idempotent on an
order-$d$ operator is
\[
 e_d(T)=\sum_{k\mid d}\mu(d/k)\frac{k}{d}
                 \sum_{j=0}^{d/k-1}T^{kj}.
\]
Support two contributes zero. In the remainder of this section $s\geq3$
and $d>2$.

For $u\in(\Z/d)^\times$, define
\[
 Q_u(b)=\frac1d\sum_{i=1}^s[ub_i]_d,
 \qquad 1\leq Q_u(b)\leq s-1,
\]
where the representatives lie in $\{1,\ldots,d-1\}$. With one fixed
character convention the Hodge dimensions are
$(Q_u-1,s-1-Q_u)$; reversing the convention reverses both entries
throughout the argument.

\begin{lemma}[Moving or all definite]\label{lem:dichotomy}
If some $u$ has $2\leq Q_u\leq s-2$, every conjugate complex projection
of connected monodromy is $\SL_h$. Otherwise every primitive embedding
is definite, and the rational monodromy of $U_c$ is finite.
\end{lemma}
\begin{proof}
Put one active branch point at infinity. There are $n=s-1$ finite
weights, and their sum is $Q_u-[ub_\infty]_d/d$. At a mixed embedding
it lies strictly between $1$ and $n-1$. This is case (a) of the good-sequence
condition in \cite[Definition 1.1]{MN}; \cite[Theorem 5.1]{MN} gives full
special-unitary closure. The residues of the finite points still have gcd
one with $d$, since the omitted residue is minus their sum.
Complexification gives $\SL_h$.

The matrices in a character basis have entries in $K_d$. Galois conjugation
transports their complex Zariski closure to that of every primitive conjugate:
a polynomial relation over $\C$ can be expanded in a finite linearly
independent list of coefficients over $K_d$. Thus fullness propagates even
to definite complex embeddings. The affine normalization and the forgetting
map to the active configuration preserve connected closure, as in
\cite[Sections 2--3]{General}.

If there is no mixed embedding, the polarization-preserving real
$K_d$-centralizer is compact. The rational primitive summand has a
monodromy-stable lattice, obtained by intersection with integral cohomology.
Its monodromy lies in a compact group and a discrete integral group, and
is therefore finite. This compactness argument and the classification of
finite cases are classical; see \cite[Theorem 8 and Lemma 11]{GV}.
We do not need that classification.
\end{proof}

\begin{lemma}[A three-point model for every all-definite factor]
\label{lem:triangle}
Suppose $Q_u(b)\in\{1,s-1\}$ for every unit $u$. Set
\[
 e=\gcd(d,b_1,b_2),\qquad m=d/e,
 \qquad \beta=(b_1/e,b_2/e,-(b_1+b_2)/e)\pmod m.
\]
Then all entries of $\beta$ are nonzero and $m\geq3$. Let $B_\beta$ be
the primitive $K_m$-factor of the Jacobian of the connected three-point
curve
\[
 y^m=x^{b_1/e}(1-x)^{b_2/e}.
\]
As rational Hodge structures, with the indicated $K_d$ action,
\[
 U_c\simeq
 \left(H^1(B_\beta,\Q)\otimes_{K_m}K_d\right)^{\oplus h}.
\]
Consequently the primitive abelian factor of $J(Y_c)$ is isogenous to
\[
 B_\beta^{\,h[K_d:K_m]}.
\]
In particular, it is an isogeny factor of a power of $J(F_N)$ whenever
$d\mid N$.
\end{lemma}
\begin{proof}
At a unit $u$ with $Q_u=1$, all $s$ positive residues sum to $d$.
Since $s\geq3$, the first two sum strictly below $d$. At a unit with
$Q_u=s-1$, apply the previous observation to $-u$: the first two original
residues then sum strictly above $d$. Neither sum is $d$, so
$b_1+b_2\not\equiv0\pmod d$. The triangle has three active entries;
its gcd with $m$ is one. A zero-sum active triangle cannot have modulus two.

The triangle's residue sum, after multiplying by a unit $u$, is $m$
when $Q_u(b)=1$ and $2m$ when $Q_u(b)=s-1$. Reduction of units from $d$
to $m$ is surjective. The rank-one $K_m$ Hodge structure
$H^1(B_\beta,\Q)$ therefore has exactly the pure types required at each
extension to $K_d$. Both sides of the displayed isomorphism have dimension
$h$ at each $K_d$ embedding and just one nonzero Hodge type there.
Any $K_d$-linear vector-space isomorphism is consequently a Hodge
isomorphism. After forgetting the $K_d$ action, scalar extension is
$[K_d:K_m]$ copies of the original rational Hodge structure.

Rational weight-one Hodge morphisms of complex abelian varieties are
homomorphisms up to isogeny. Thus the preceding comparison supplies actual
algebraic maps on the whole primitive factor, not only matching Hodge
numbers. Finally, on $F_m$ in affine coordinates $u^m+v^m=1$, the map
$x=u^m$, $y=u^{b_1/e}v^{b_2/e}$ dominates the triangle curve. Pullback and
norm make $B_\beta$ an isogeny factor of $J(F_m)$, and the power map
$F_N\to F_m$ finishes the claim.
\end{proof}

\begin{remark}
The lemma avoids a conductor-by-conductor list and any degeneration or
specialization theorem. It includes the quartic word $(1,1,1,1)$, the
sextic word $(1,1,1,1,1,1)$, and the conductor-fifteen word $(1,2,4,8)$.
Its conclusion is about the primitive factor, not the whole cyclic
Jacobian. Its exact historical priority has not been established.
\end{remark}

\section{Simultaneous moving monodromy at every exponent}
Retain only the primitive orbits with a mixed embedding. For an individual
complex character write $q_i=\zeta_N^{c_i}$ at its active positions; their
product is one. Quadratic characters are considered separately below.

\begin{theorem}[Exact moving blocks]\label{thm:blocks}
For nonquadratic moving characters, different active supports belong to
independent simple factors. On a fixed support of size at least five,
the blocks are $\{c,-c\}$, with standard and dual $\SL_{s-2}$
occurrences. On four active labels, the blocks are
\[
 (\mathcal V_4 c\cup-\mathcal V_4 c)\cap C,
\]
where $\mathcal V_4$ is the Klein four group and the intersection is with
the actual character code. The oriented $\mathcal V_4$ identifications
are algebraic curve-graph correspondences; negation is supplied by
polarization with its dual and Tate twist.

Every nonzero quadratic factor has full $\Sp_{s-2}$ monodromy.
Different quadratic supports are independent, and no quadratic moving
factor is joined to a nonquadratic one. The connected complex linear
monodromy is the product of these groups with the stated occurrences;
all-definite factors have trivial connected action.
\end{theorem}

\subsection{Common projective representations and active support}
All projective representations are defined on the same $\pi_1(M_{0,r})$.
Choices of lifts change a character operator only by a deck scalar. The
normalizer of a graph in $H\times H$, for a centerless simple group $H$,
is that graph: a normalizing pair $(a,b)$ has $\theta(a)^{-1}b$ central.
Since the connected closure is normal in the full closure, any connected
graph relation therefore constrains the original braid images, before
finite base changes take powers of their eigenvalues.

An active pair twist is a nonzero rank-one modification of the identity.
This follows from \cite[Theorems 2.2 and 2.5]{MN}: put a different active
point at infinity; the spanning vectors have their sole relation
$\sum_i g_i=0$, the chosen vector is nonzero, and the Hermitian form and
the rank-one coefficient are nonzero. Its spectrum is
\[
 (q_iq_j,1,\ldots,1).
\]
When $q_iq_j=1$ it is a \emph{nonidentity} unipotent. Forgetting an inactive
label makes that same twist projectively trivial. This separates unequal
supports whenever the simple target groups are isomorphic; otherwise
independence is automatic. Goursat's lemma then reduces the joint
classification to graph relations on a single support.

\subsection{Rank at least three: the existing parity-free argument}
This step is the argument of \cite[proof of Proposition 6.2]{MN}, applied
to two possibly different residue words rather than two embeddings of one
word. It is not a new method. After choosing inner or inverse-transpose
orientation, a projective graph gives one simultaneous conjugating matrix.
For an inner graph, comparison of the pair spectra above forces the scalar
multipliers to be one: the eigenvalue of multiplicity $h-1\geq2$ is
canonically one, including the unipotent case. Thus the canonical affine
pure-braid representations are \emph{linearly} conjugate on every pair
generator, hence on the entire pure braid group.

For a twist enclosing three finite active points $i,j,k$,
\cite[Theorem 2.8]{MN} gives the trace
\[
 (h-2)+2q_iq_jq_k.
\]
The pair products and these triple products agree for the two words.
Their quotient recovers each $q_i$, with no division by two in $\Z/N$.
The missing infinity root follows from the product-one relation. In the
dual orientation replace all roots by their inverses. This proves that
the words are $c$ or $-c$.

It is important that the scalar normalization is simultaneous on generators.
Comparing an isolated projective triple spectrum is not a substitute for
linear conjugacy and does not justify this step.

\subsection{Rank two: the remaining half-turn is detected by a trace}
Put $u=q_1$, $v=q_2$, $w=q_3$, so $q_4=(uvw)^{-1}$.
In the basis $(g_1,g_2)$ furnished by the source Gram and pair formulas,
the adjacent pair twists have matrices
\begin{equation}\label{eq:AB}
 A=\begin{pmatrix}uv&v(1-u)\\0&1\end{pmatrix},\qquad
 B=\begin{pmatrix}1&0\\1-w&vw\end{pmatrix}.
\end{equation}
For example, the upper-right coefficient of $A$ is obtained by substituting
$\langle g_2,g_1\rangle$ into the rank-one formula; the lower-left
coefficient of $B$ is obtained from $\langle g_1,g_2\rangle$.
Thus the same source conventions determine both operators. Direct
multiplication gives
\[
 \det(XI-AB)=(X-v)(X-uvw).
\]
The scalar-conjugacy invariant
\begin{equation}\label{eq:J}
 \mathcal J(A,B)=
 \frac{\tr(A)\tr(B)\tr(AB)}{\det A\det B}
 =\frac{(1+uv)(1+uw)(1+vw)}{uvw}
 =\sum_{i=1}^4(q_i+q_i^{-1})
\end{equation}
is unchanged by independent scalar changes of $A,B$ and by simultaneous
conjugation. It is a classical trace-type invariant; no invention of a new
character-variety invariant is claimed.

Pair spectra determine each of
$X=q_1q_2$, $Y=q_1q_3$, $Z=q_1q_4$ up to inversion. The four even
inversions are realized by $\mathcal V_4$; simultaneous inversion of
all roots supplies the other four. After applying one of these known
geometric or dual transformations, the two words have equal pair
products. Their coordinate ratios are then one common number $\epsilon$
with $\epsilon^2=1$. The only extra candidate is
\[
 (q_1,q_2,q_3,q_4)\longmapsto(-q_1,-q_2,-q_3,-q_4).
\]
It corresponds to adding $N/2$ to every active residue when $N$ is even.
If it produces an inactive entry, support separation already excludes it.

Under this half-turn, $\mathcal J$ changes sign. Hence a projective graph
is impossible when $\mathcal J\neq0$. If $\mathcal J=0$, one pair product
is $-1$, and the complementary pair product is also $-1$. Let $P$ swap
the entries within these two pairs. Then
\[
 -q_i=q_{P(i)}^{-1}\quad\text{for every }i.
\]
The half-turn is therefore already in $-\mathcal V_4$ and supplies no
additional graph. This proves the rank-two necessity for every modulus.
No coprime pair-sum hypothesis is needed.

For the converse, a double transposition of four branch points is induced
by a M\"obius transformation preserving their cross-ratio. For the two
cyclic equations the quotient of their pulled-back Kummer functions has
divisor divisible by the actual conductor. Since $\operatorname{Pic}^0
(\PP^1)=0$, adjoining a root of a parameter-dependent constant gives a
deck-equivariant isomorphism after a finite base cover. Quotient pullback
and norm transfer its graph to the primitive character factors. This is
exactly the oriented graph construction in \cite[Section 3]{General},
which does not require odd conductor. Polarization realizes the negative
orientation; it is not an unsupported degree-zero Hodge isomorphism to
the opposite character.

\paragraph{An even-order regression.}
At modulus eight, $c=(1,1,1,5)$ and $e=(5,5,5,1)$ have identical
projective pair spectra but are not in one $\pm\mathcal V_4$ orbit.
Their trace invariants are $2\sqrt2$ and $-2\sqrt2$. Both primitive
orbits have mixed embeddings. Thus pair spectra alone would permit a
false graph; the trace calculation excludes it.

\subsection{Quadratic factors and completion of the product argument}
A quadratic character is determined by its even active support. Its
quotient is hyperelliptic and has full symplectic connected monodromy,
by A'Campo's theorem with the compact homology quotient; the ordered and
normalized configurations preserve connected closure. The explicit source
match is in \cite[Section 7]{General}.

Distinct quadratic supports separate by a nonzero Picard--Lefschetz
transvection on an active pair, trivial after forgetting an inactive label.
The only isomorphism between a symplectic and a special-linear simple
factor in the dimensions at issue is $\Sp_2=\SL_2$. On a common
four-point support the quadratic pair twists are all unipotent. A graph
to a nonquadratic rank-two factor would force $q_iq_j=1$ for every pair,
so all active roots would be $-1$. That is itself quadratic, a contradiction.

The connected solvable radical acts trivially on every full simple
projection and on every finite factor, hence is trivial. Semisimple
Goursat theory now assembles the pairwise classification; there are no
additional relations invisible in all pairwise projections. Lifting from
adjoint groups gives the stated product of simply connected $\SL$ and
$\Sp$ groups. Every factor has a faithful standard occurrence in $H^1$,
so this product is the connected linear image. This completes
Theorem~\ref{thm:blocks}.

\section{Algebraic sources and the Hodge reduction}
\subsection{The classical determinant correspondence has no parity restriction}
For a degree-$d>2$ primitive cyclic factor of rank $h$, put
\[
 D_c\otimes\C=\bigoplus_{u\in(\Z/d)^\times}\bigwedge^h V_{uc}.
\]
The required input is the rational algebraic realization of $D_c$ as a
summand of $H^h(F_d^h,\Q)$. This is the classical Schoen determinant
construction \cite[Lemmas 1.1--1.2]{Schoen88}; its finite-correspondence
normalization is written out in \cite[Section 7]{Structural} and in
\cite{Fermat}. Here is the specific parity-independent mechanism.

On $Z=Y_c^h$, average over
$H=\ker(\mu_d^h\to\mu_d)\rtimes\mathfrak S_h$ and apply the primitive
cyclotomic projector for a deck action in \emph{one} coordinate. Subtracting
the full $\mu_d^h\rtimes\mathfrak S_h$ average removes the invariant
part. The unsigned geometric permutation average is the ordinary exterior
antisymmetrizer on degree-one tensors because of the Koszul sign. The
projector's cohomological image is exactly $D_c$.

The normal quotient $Z/H$ is a cyclic cover of $\Sym^h\PP^1=\PP^h$
branched along $h+2$ evaluation hyperplanes. They are in general position,
and taking $d$th roots gives a diagonal Fermat $h$-fold mapping finitely
onto the same quotient. The fiber-product correspondence, its transpose,
and division by their finite degrees give inverse maps on the projected
summand. These are proper algebraic correspondences; their compositions
are the finite deck averages. Neither root-taking, averaging, nor the
primitive idempotent assumes $d$ odd.

Using the algebraic $H^1$ correspondence with the Jacobian gives the same
summand in its cohomology. Basepoint changes in Abel--Jacobi are translations
and act trivially on $H^1$. Acting in every coordinate simultaneously would
instead multiply the character by $h$, and can lose it when $d\mid h$;
that operation is not used. No new cycle-construction priority is claimed.

\subsection{Rational invariant spanning, including finite factors}
The first fundamental theorems for $\SL$ and $\Sp$ express invariant
tensors using contractions and, for $\SL$, full volumes. Apply them to
Theorem~\ref{thm:blocks}. Actual oriented graphs identify occurrences,
polarization handles duals, and alternation passes from tensor powers of
$H^1$ to $H^k(A_b^q)$, with nonzero rational factorial normalizations.
Quadratic symplectic blocks require only polarization contractions.

Use \emph{full rational} tensor products of the $D_c$, the
$H^1(B_\beta,\Q)$ furnished by Lemma~\ref{lem:triangle}, and Tate factors.
Support-three factors are included in the same lemma. The finite factors'
whole first cohomology, not merely their determinants, is inserted by the
isogenies proved there. All forward maps are algebraic.

For rational spanning, temporarily extend scalars to $L=\Q(\zeta_N)$.
Every character selector on a raw slot is
\[
 e_c=\frac1{|G|}\sum_{g\in G}\zeta_N^{-c(g)}g^*,
\]
an $L$-linear combination of rational algebraic deck correspondences.
Expand the selectors in each complex invariant diagram, including
selectors on contractions and graph maps. Each diagram is consequently
in the $L$-span of images of rational algebraic maps from the full sources.
Each expanded map still has invariant image: it uses determinant volumes,
finite-factor slots, equivariant deck or graph maps, and polarization
contractions. In a fixed degree only finitely many diagrams and terms
are required. Their direct sum is a rational Hodge surjection onto the
invariant space, since it is surjective after extension to $L$.

This is the full-rational-source argument of
\cite[Section 7]{Structural}, now with Lemma~\ref{lem:triangle} replacing
its special conductor-fifteen source and with quadratic pairings added.
It assumes neither rank one over a common cyclotomic field nor independence
of different cyclotomic CM characters. In particular it controls mixed
Hodge classes rather than inferring them from factorwise HC. This proves
Theorem~\ref{thm:presentation}.

\subsection{From one Fermat Jacobian to all source products}
Assume $\HC(J(F_N)^q)$ for all $q$. For $N\geq3$, Abel--Jacobi pullback
surjects on the cohomology of $F_N$, hence on all its powers. Semisimplicity
lifts Hodge classes through this surjection, proving HC for curve products.
For $N=2$ these products are rational and the conclusion is immediate.

The Shioda--Katsura rational map
$F_N^{h-1}\times F_N^1\dashrightarrow F_N^h$ of degree $N$ is resolved
by blowing up $F_N^{h-2}\times F_N^0$ \cite[Lemmas 1.1--1.2]{SK}.
The blowup formula and $N^{-1}f_*f^*=1$ reduce HC for arbitrary products
of Fermat hypersurfaces inductively to curve products and points. Both
the source and the blowup center are controlled; arbitrary rational
domination alone would not justify the deduction. For every $d\mid N$,
the finite power map $F_N^h\to F_d^h$ gives the same conclusion for
products of degrees dividing $N$. Effective Tate factors are supplied by
projective spaces. Thus all Hodge classes in every source of
Theorem~\ref{thm:presentation} have algebraic representatives in their
specified realizations.

At a very general fiber all rational Hodge classes on all powers are
fixed by connected monodromy. The generic-monodromy and algebraic
Hodge-locus argument, including the countable union across powers and
degrees, is recorded with source locators in \cite[Section 6]{General}.
Lift any such class through the rational Hodge surjection of
Theorem~\ref{thm:presentation}. Polarizable rational Hodge structures
are semisimple, so the lift can be chosen Hodge. Its algebraic source
representatives map forward to an algebraic representative on $A_b^q$.
Only the forward maps must be algebraic; an algebraic splitting of the
surjection is not assumed. This proves the forward implication of
Theorem~\ref{thm:transfer}.

For the reverse implication take $G=(\Z/N)^2$ and branch vectors
$(1,0),(0,1),(-1,-1)$. Its three-point cover is the Fermat curve
$x^N+y^N=1$, via the coordinate $x^N$. Its full branch space is a point.
The right-hand assertion therefore includes HC for all powers of $J(F_N)$.
Corollary~\ref{cor:aoki} follows from the established Fermat degrees.

For Corollary~\ref{cor:category}, the same Shioda--Katsura blowup induction,
without any Hodge-conjecture premise, puts the cohomology of every relevant
Fermat product in $\langle H^1(F_N),\Q(1)\rangle$. Semisimplicity and
Theorem~\ref{thm:presentation} put every invariant space in that category.
The reverse inclusion again uses the rigid Fermat datum. Thus this
categorical fixed-part description is independent of the cycle-algebraicity
premise, although its presentation is by the algebraic maps already proved.


\section{Verification boundary and novelty}
The accompanying exact-arithmetic script verifies five symbolic matrix
identities, the four-point reconstruction for all 236,096 active zero-sum
words at moduli 2 through 32, and separation of 46,032 nongraph half-turn
candidates. It also verifies the common pair-kernel and triple reconstruction,
and the triangle types for 391 primitive all-definite four-point words with
first-row residue sum equal to the conductor, at conductors through 32.
Negation accounts for the other definite first-row sum. Thirty-six tested
words have a proper triangle conductor. Higher-support examples are stated
explicitly rather than included in that four-point count.

These are symbolic and bounded arithmetic checks, not proofs of geometric
monodromy, algebraic correspondences, or the Hodge conjecture. The universal
arguments are the written proofs above. No new Lean theorem is asserted;
the repository-wide Lean and legacy certificate suites were not run in
this continuation environment. None of the old frozen sources was changed.

The credible contribution under examination is the exact simultaneous
classification, with its all-exponent Fermat-source assembly for arbitrary
noncyclic deck groups. Individual density, the high-rank trace method,
finite-monodromy classification, determinant cycles, and the established
Fermat degrees are earlier work. The rank-two trace application and the
three-point model are short arguments, not evidence by themselves of a
major new mechanism. The complete 1989 Schoen chapter and the full
hypergeometric-equivalence literature remain priority checks; a failed
search is not a proof of novelty. Neither Duke nor JEMS readiness is
asserted by this draft.

\begin{thebibliography}{99}
\bibitem{MN} A. Menet and D.-M. Nguyen,
\emph{Representations of braid groups via cyclic covers of the sphere:
Zariski closure and arithmeticity}, arXiv:2310.10401v3.
Definitions and locators: Definition 1.1; Theorems 2.1, 2.2, 2.5, 2.8,
5.1; Proposition 6.2 and Remark 6.4.
\bibitem{GV} E. Ghate and T. N. Venkataramana,
\emph{Sums of fractions and finiteness of monodromy},
arXiv:1612.04548; Indagationes Mathematicae (2017),
\href{https://doi.org/10.1016/j.indag.2017.09.003}{doi:10.1016/j.indag.2017.09.003}.
Theorem 8 and Lemma 11.
\bibitem{Schoen88} C. Schoen, \emph{Hodge classes on self-products of a
variety with an automorphism}, Compositio Mathematica 65 (1988), 3--32.
Lemmas 1.1--1.2 and Remark 1.4.
\bibitem{Schoen89} C. Schoen, \emph{Cyclic covers of $\PP^v$ branched
along $v+2$ hyperplanes and the generalized Hodge conjecture for certain
abelian varieties}, Lecture Notes in Mathematics 1399 (1989), 137--154.
\href{https://doi.org/10.1007/BFb0095974}{doi:10.1007/BFb0095974}.
Full chapter not re-read in this continuation; no new exact theorem
attribution is made to it.
\bibitem{SK} T. Shioda and T. Katsura,
\emph{On Fermat varieties}, Tohoku Mathematical Journal 31 (1979), 97--115.
\bibitem{Aoki00} N. Aoki, \emph{Some remarks on the Hodge conjecture for
abelian varieties of Fermat type}, Commentarii Mathematici Universitatis
Sancti Pauli 49 (2000), 177--194, Theorem 0.1.
\bibitem{Aoki02} N. Aoki, \emph{Hodge cycles on CM abelian varieties of
Fermat type}, Commentarii Mathematici Universitatis Sancti Pauli 51 (2002),
99--130, Theorem 1.2. This author's restatement was read in the available
full-text mirror; that mirror's landing-page title is mismatched.
\bibitem{General} K. Kistner, \emph{All powers of very general Jacobians
of elementary abelian covers}, repository manuscript
\texttt{manuscripts/general\_all\_powers.tex}, checkpoint specified above.
Written project proof, not an externally refereed input.
\bibitem{Structural} K. Kistner, repository structural companion,
\texttt{manuscripts/structural\_extensions.tex}, same checkpoint,
Section 7. Written project proof, not an externally refereed input.
\bibitem{Fermat} K. Kistner, repository Fermat-transfer companion,
\path{manuscripts/fermat_determinant_transfer.tex}, same checkpoint.
Explicit presentation of a classical mechanism, not a priority claim.
\end{thebibliography}
\end{document}
